% =============================================================================
% Question 3 - controlled-condition detection
% =============================================================================
\clearpage
\hypertarget{q3-functional-blocks}{}
\sectiontitle{Q3A. Selection of the Critical Functional Blocks}

\evidencebox{Concept-exploration question:}{\textbf{Can the integrated slab
detect controlled conditions representing slipping, perspiration, thermal
buildup, and potentially harmful localized loading?}}

{\fontsize{9.35}{11.05}\selectfont
A functional block was retained when it (1) transforms a physical input into a
distinct and testable output, (2) contains uncertainty that is not resolved by
the literature for this directly printed 5.2 mm PDMS slab, (3) can fail
independently of the other blocks, and (4) would invalidate a downstream claim
if it failed. This method keeps the decomposition tied to technical feasibility
rather than to arbitrary component names. Published solutions are leveraged
where they already establish a mechanism; project-specific proof is reserved
for transfer to the present geometry, material stack, fabrication process, and
coordinate system \cite{albrecht2019,liew2020,aitmammar2020,q3kwak2020}.\par}

\vspace{0.45em}
\begingroup
\fontsize{7.8}{9.15}\selectfont
\setlength{\tabcolsep}{2.5pt}
\renewcommand{\arraystretch}{1.12}
\begin{tabularx}{\linewidth}{>{\centering\arraybackslash\bfseries\color{AUBRed}}p{0.72cm}
                            >{\raggedright\arraybackslash\bfseries}p{3.25cm}
                            Y
                            >{\raggedright\arraybackslash}p{4.45cm}}
\rowcolor{AUBRed}
\color{white}\textbf{ID} & \color{white}\textbf{Functional block} &
\color{white}\textbf{Why it is isolated} &
\color{white}\textbf{First technical-feasibility question} \\
FB1 & Interface mechanics and contact transfer & Every sensor depends on how
the compliant slab transfers concentrated and distributed loading. & Does the
5.2 mm PDMS geometry preserve a usable local response without unacceptable
spatial spreading or mechanical cross-coupling? \\
\rowcolor{AUBPale}
FB2 & Normal-pressure and directional-shear transduction & Force magnitude and
direction arise from the same deformable interface but require separable
outputs \cite{albrecht2019}. & Can normal and two-axis tangential loading be
recovered over the selected range after integration? \\
FB3 & Temperature transduction & Printed silver resistance is temperature
sensitive, but bending, strain, encapsulation, and thermal mass can alter the
response \cite{liew2020,q3kwak2020}. & Does the embedded element measure local
temperature with adequate accuracy, response, and recovery? \\
\rowcolor{AUBPale}
FB4 & Vapour humidity and liquid-sweat interface & Vapour access is required,
whereas direct liquid exposure can saturate or damage the sensing route
\cite{aitmammar2020,paterno2020}. & Can RH change and liquid exposure be
distinguished while the protected route remains recoverable? \\
FB5 & Printed-PDMS circuit integrity and encapsulation & Wetting, swelling,
adhesion, cracking, leakage, and cyclic strain are material/process risks
\cite{albrecht2018stretch,sun2016}. & Do the printed conductors and insulation
remain continuous and stable after environmental and mechanical exposure? \\
\rowcolor{AUBPale}
FB6 & Signal conditioning, calibration, and synchronization & Multimodal
channels require independent calibration and a common time base before their
outputs can be compared \cite{q3tang2022,q3kwak2020}. & Can raw channels be
converted into traceable physical units without timing or cross-talk errors? \\
FB7 & Slip and relative-motion discrimination & Shear is a force condition;
slip is relative motion and therefore needs an independent displacement
reference \cite{noll2018,q3tang2022}. & Can true translation be distinguished
from high shear without motion? \\
\rowcolor{AUBPale}
FB8 & Coordinate registration and multimodal interpretation & Measurements
become useful only when mapped to the correct slab location while retaining
modality-specific evidence \cite{q3turner2021}. & Can event type and location
be recovered without misleading interpolation or averaging? \\
\end{tabularx}
\endgroup

\sourceanchor{Functional blocks selected from the concept's physical signal
path and from unresolved transfer risks identified in the sensor, prosthetic
interface, and printed-PDMS literature
\cite{albrecht2019,liew2020,aitmammar2020,albrecht2018stretch,
q3tang2022,noll2018,q3turner2021}.}

\vspace{0.55em}
\evidencebox{Selection result:}{Eight blocks are retained, exceeding the six
required by the assignment. FB1--FB6 establish whether the slab can produce
traceable measurements; FB7--FB8 establish whether those measurements support
the stated condition and coordinate. A failure in an upstream block is not
averaged into a downstream software score.}

\clearpage
\sectiontitle{Q3B. Definition and Boundaries of the Functional Blocks}

{\fontsize{9.35}{11.05}\selectfont
Each block below has an explicit input, transformation, and output and is tied
to the engineering discipline that governs its dominant uncertainty. The
boundaries prevent one successful subsystem from being used as evidence for a
different subsystem. The validation fixtures and reference instruments are
not product blocks; they are the independent means used later to examine the
eight product blocks.\par}

\vspace{0.45em}
\begingroup
\fontsize{7.75}{9.1}\selectfont
\setlength{\tabcolsep}{2.5pt}
\renewcommand{\arraystretch}{1.12}
\begin{tabularx}{\linewidth}{>{\raggedright\arraybackslash\bfseries\color{AUBRed}}p{2.75cm}
                            >{\raggedright\arraybackslash}p{5.55cm}
                            >{\raggedright\arraybackslash}p{3.15cm}Y}
\rowcolor{AUBRed}
\color{white}\textbf{Block} &
\color{white}\textbf{Input $\rightarrow$ transformation $\rightarrow$ output} &
\color{white}\textbf{Dominant discipline} &
\color{white}\textbf{Boundary of responsibility} \\
FB1: interface transfer & Contact force, heat, and moisture $\rightarrow$ PDMS
deformation/diffusion $\rightarrow$ condition delivered at an active region. &
Mechanical engineering; transport phenomena. & Ends at the local stimulus
reaching the sensor plane; it does not include electrical conversion. \\
\rowcolor{AUBPale}
FB2: pressure/shear & Normal and tangential stress $\rightarrow$ four-capacitor
geometry $\rightarrow$ calibrated $F_z,F_x,F_y$ or stresses. & Mechanical and
electrical engineering. & Includes axis separation and mechanical cross-talk;
excludes deciding whether motion occurred. \\
FB3: temperature & Local heat flux/temperature $\rightarrow$ resistance change
$\rightarrow$ calibrated temperature trace. & Thermal and electrical
engineering. & Includes thermal lag and strain sensitivity; excludes clinical
interpretation of temperature. \\
\rowcolor{AUBPale}
FB4: humidity/sweat route & Water vapour or liquid exposure $\rightarrow$
protected dielectric response $\rightarrow$ RH-related signal plus liquid-state
flag. & Materials science; electrical engineering. & Includes vapour access,
saturation, drying, and protection; RH is not treated as sweat rate. \\
FB5: printed hardware & Printing, curing, bending, wetting, and cycling
$\rightarrow$ conductor/insulator state $\rightarrow$ continuity, leakage,
adhesion, and baseline stability. & Materials science; manufacturing. & Covers
the physical circuit stack, not sensor calibration or software compensation. \\
\rowcolor{AUBPale}
FB6: measurement chain & Raw capacitance/resistance and timestamps $\rightarrow$
conditioning, calibration, filtering, and synchronization $\rightarrow$
physical-unit channels with quality flags. & Electrical, instrumentation, and
data engineering. & Ends before event classification; retains raw data and
calibration provenance. \\
FB7: slip discrimination & Synchronized pressure/shear patterns plus an
independent displacement reference $\rightarrow$ event comparison
$\rightarrow$ slip/no-slip, direction, and onset error. & Mechatronics and
signal processing. & Requires motion evidence; high shear alone remains a
force event. \\
\rowcolor{AUBPale}
FB8: spatial interpretation & Calibrated channels plus CAD coordinates
$\rightarrow$ registration/localization $\rightarrow$ modality-specific event
location and bounded interface-condition display. & Computer science/software
engineering. & Includes coverage and confidence; excludes a validated clinical
discomfort or injury score. \\
\end{tabularx}
\endgroup

\vspace{0.55em}
\begin{center}
\begin{tikzpicture}[>=stealth,node distance=4mm]
  \tikzstyle{fbbox}=[draw=AUBLine,fill=AUBPale,rounded corners=1pt,
    minimum height=10mm,text width=32mm,align=center,
    font=\fontsize{7}{8}\selectfont]
  \node[fbbox] (input) {controlled interface\\condition};
  \node[fbbox,right=of input] (physical) {FB1--FB5\\physical transfer\\and sensors};
  \node[fbbox,right=of physical] (chain) {FB6\\traceable\\measurement chain};
  \node[fbbox,right=of chain] (meaning) {FB7--FB8\\event and coordinate};
  \draw[->,thick,color=AUBRed] (input)--(physical);
  \draw[->,thick,color=AUBRed] (physical)--(chain);
  \draw[->,thick,color=AUBRed] (chain)--(meaning);
  \draw[->,thick,dashed,color=AUBGray]
    ([yshift=-4mm]meaning.south west) to[bend left=20]
    node[below,font=\fontsize{6.8}{7.8}\selectfont]{independent reference closes the validation loop}
    ([yshift=-4mm]input.south east);
\end{tikzpicture}
\end{center}
{\fontsize{7.55}{8.75}\selectfont\color{AUBGray}
\textbf{Figure Q3-0A. Functional-block signal path.} Physical transduction,
measurement conversion, and interpretation remain separate claim layers. The
external reference evaluates the chain but is not itself part of the product.\par}

\clearpage
\sectiontitle{Q3C. Literature Evidence and Project-Specific Proof}

{\fontsize{9.25}{10.95}\selectfont
The literature establishes that the individual sensing and interpretation
mechanisms are plausible, but no reviewed study implements this exact set of
directly printed sensors in the present 5.2 mm slab. The matrix therefore
separates \textbf{knowledge that can be leveraged} from the limited set of
\textbf{project-specific relationships that must be established}. The detailed
literature plots, fixtures, equations, protocols, and decision metrics in the
remainder of Q3 provide the evidence route for these relationships.\par}

\vspace{0.4em}
\begingroup
\fontsize{7.55}{8.9}\selectfont
\setlength{\tabcolsep}{2.35pt}
\renewcommand{\arraystretch}{1.1}
\begin{tabularx}{\linewidth}{>{\raggedright\arraybackslash\bfseries\color{AUBRed}}p{2.5cm}
                            >{\raggedright\arraybackslash}p{6.15cm}Y}
\rowcolor{AUBRed}
\color{white}\textbf{Functional block} &
\color{white}\textbf{Known and leveraged from prior work} &
\color{white}\textbf{Relationship established for this prototype} \\
FB1: interface transfer & Prosthetic-interface loading is cyclic,
multidirectional, location-dependent, and sensitive to contact geometry; PDMS
also introduces viscoelastic recovery \cite{q3sanders1992,q3tang2022,
albrecht2019}. & Transfer from a known contact area and coordinate through the
5.2 mm slab, including spatial spread, recovery, and coupling between adjacent
regions. \\
\rowcolor{AUBPale}
FB2: pressure/shear & A four-capacitor printed cell can separate common-mode
normal loading from directional shear over a demonstrated force range, while
also exhibiting priming, drift, and hysteresis \cite{albrecht2019}. & The
specimen-specific sensitivity matrix, usable range, axis error, hysteresis, and
cross-sensitivity after printing and integration. \\
FB3: temperature & Printed silver meanders show an approximately linear
resistance--temperature relationship, and soft prosthetic sensors can track
skin-interface temperature \cite{liew2020,q3kwak2020}. & Calibration near the
skin-relevant range, thermal response/recovery, spatial leakage, and separation
of temperature from force- and strain-induced resistance changes. \\
\rowcolor{AUBPale}
FB4: humidity/sweat route & Printed capacitive RH sensors respond to vapour;
prosthetic-liner RH can continue rising after temperature begins to recover
\cite{aitmammar2020,paterno2020}. & Vapour response through the selected
protection, liquid-sweat saturation and drying behaviour, temperature coupling,
and post-exposure drift or damage. \\
FB5: printed hardware & Oxygen-plasma treatment enables AgNP printing on PDMS;
solvent absorption can swell PDMS, and partially embedded conductors improve
mechanical retention \cite{albrecht2018stretch,sun2016}. & Continuity,
insulation, adhesion, resistance, cracking, and calibration stability for the
actual ink--printer--PDMS process under bending, heat, moisture, and cycling. \\
\rowcolor{AUBPale}
FB6: measurement chain & Published systems synchronize stress with motion and
multiple skin-interface modalities while retaining independent signals
\cite{q3tang2022,q3kwak2020}. & Channel calibration, common-clock accuracy,
sampling and packet-loss behaviour, filtering delay, and traceability from raw
channel to physical unit. \\
FB7: slip discrimination & Relative limb--socket motion can be measured against
an independent displacement reference, and stress and motion are related but
not interchangeable \cite{noll2018,q3tang2022}. & Separation of translated
contact from high-shear/no-motion trials, including direction, onset time,
preload dependence, and forward/reverse consistency. \\
\rowcolor{AUBPale}
FB8: spatial interpretation & Socket pressure data can be registered to sensor
coordinates and displayed using several interpolation methods
\cite{q3turner2021}. & Correct sensor-to-CAD mapping, localization error at
sensor centres and gaps, validated coverage, combined-event separation, and a
display that preserves disagreement between modalities. \\
\end{tabularx}
\endgroup

\sourceanchor{Known mechanisms and transferable solutions are drawn from the
printed-sensor, prosthetic-interface, multimodal monitoring, relative-motion,
and spatial-visualization literature. Only the right-hand column is specific to
the present prototype \cite{albrecht2019,liew2020,aitmammar2020,
albrecht2018stretch,sun2016,q3kwak2020,noll2018,q3turner2021}.}

\vspace{0.55em}
\evidencebox{Risk-order consequence:}{FB1--FB6 are examined before FB7--FB8.
This follows the assignment's feasibility logic: an event classifier or smooth
map cannot rescue inaccurate transduction, damaged conductors, unsynchronized
channels, or an unprotected humidity route.}

\clearpage
\sectiontitle{Q3C. Operational Definitions and Claim Boundary}

{\fontsize{9.55}{11.3}\selectfont
Q3 is a \textbf{bench-validation question}. A known physical condition is
applied at a known time and coordinate, and the integrated 5.2 mm PDMS slab is
evaluated for the correct sensing modality, direction where applicable,
location, magnitude, and duration. This stage addresses bench performance;
participant discomfort and skin-injury thresholds belong to a subsequent
clinical-validation stage. Accordingly, the admissible claim is detection of
\emph{simulated interface-condition events}; a clinical discomfort-risk claim
requires comparison with user-reported discomfort, skin observations, and
anatomically justified thresholds \cite{q3armitage2026}.\par}

\hypertarget{q3-operational-definitions}{}
\qtwosubsection{3.1 Operational meaning of each condition}

\begingroup
\fontsize{8.55}{10.05}\selectfont
\setlength{\tabcolsep}{3.2pt}
\renewcommand{\arraystretch}{1.14}
\begin{tabularx}{\linewidth}{>{\raggedright\arraybackslash\bfseries\color{AUBRed}}p{2.6cm}
                            >{\raggedright\arraybackslash}p{4.0cm}Y Y}
\rowcolor{AUBRed}
\color{white}\textbf{Reported condition} &
\color{white}\textbf{Controlled ground truth} &
\color{white}\textbf{Required evidence from the slab} &
\color{white}\textbf{Evidence boundary at this stage} \\
Localized loading & Known normal force, contact area, coordinate, and loading
time. & Pressure channel responds at the correct coordinate; magnitude follows
calibration; neighbouring and non-pressure channels remain bounded. & The load
is painful, unsafe, or tissue-damaging. \\
\rowcolor{AUBPale}
Directional shear & Known tangential force in $+x,-x,+y,-y$ under a documented
normal preload. & Correct shear direction and magnitude are recovered without
being confused with normal compression. & A shear peak by itself proves that
the liner slipped. \\
Slipping or pistoning & Known relative translation measured independently by a
stage, displacement gauge, encoder, or video reference. & A time-linked change
in shear, pressure distribution, or detected coordinate agrees with the
reference displacement. & Slip is inferred from one high shear reading without
independent motion evidence. \\
\rowcolor{AUBPale}
Thermal buildup & Known temperature step, ramp, or sustained heated contact at
a known coordinate. & Temperature follows the reference with measured error,
response time, recovery time, and limited mechanical cross-sensitivity. & The
temperature constitutes a clinical burn or injury threshold. \\
Perspiration-related condition & Vapour exposure expressed as relative humidity
(RH), tested separately from a known liquid artificial-sweat dose. & The
humidity channel detects vapour; liquid exposure, saturation, drying, recovery,
and permanent drift are separately reported. & RH is treated as sweat rate, or
liquid presence is assumed from RH alone. \\
\end{tabularx}
\endgroup
\sourceanchor{Operational definitions used for the current validation, anchored to published
prosthetic-interface pressure, shear, temperature, and microclimate evidence
\cite{q3sanders1992,q3tang2022,paterno2020,q3kwak2020}.}

\vspace{0.2em}
\begin{center}
\begin{tikzpicture}[>=stealth,node distance=3mm]
  \tikzstyle{opbox}=[draw=AUBLine,fill=AUBPale,rounded corners=1pt,
    minimum height=7mm,text width=34mm,align=center,
    font=\fontsize{7}{8}\selectfont]
  \node[opbox] (g) {independent\\ground truth};
  \node[opbox,right=of g] (s) {slab modality\\and coordinate};
  \node[opbox,right=of s] (c) {reference\\comparison};
  \node[opbox,right=of c] (b) {bounded engineering\\claim};
  \draw[->,thick,color=AUBRed] (g)--(s);
  \draw[->,thick,color=AUBRed] (s)--(c);
  \draw[->,thick,color=AUBRed] (c)--(b);
\end{tikzpicture}
\end{center}
{\fontsize{7.5}{8.7}\selectfont\color{AUBGray}
\textbf{Figure Q3-0. Claim-control chain.} The sensor output cannot label its
own test; every event requires an external reference and a limited claim.\par}

\vspace{0.55em}
\evidencebox{Central distinction:}{\textbf{Shear is a force condition; slip is a
motion condition.} Slip can generate shear, but slip detection requires a
displacement ground truth. Likewise, humidity describes water vapour, whereas
perspiration can also produce liquid sweat. These conditions are therefore
labelled and tested separately.}

\clearpage
\sectiontitle{Q3C. Quantitative Literature Evidence}

\hypertarget{q3-mechanical-evidence}{}
\qtwosubsection{3.2 Mechanical loading and motion}

{\fontsize{9.55}{11.3}\selectfont
Simultaneous pressure and shear measurement is necessary because interface
loading is multidirectional. Sanders \textit{et al.} measured resultant shear
of 5.6--39.0 kPa at the tibial flares and 5.0--40.7 kPa posteriorly during
walking; the waveforms were commonly double-peaked, with peaks occurring at
25--40\% and 65--85\% of stance \cite{q3sanders1992}. Tang \textit{et al.}
combined multidirectional stress sensing with three-dimensional motion capture
and reported posterior-proximal peak pressure of 55--59 kPa and shear of
12--19 kPa across nine sessions over 12 months in one transfemoral participant
\cite{q3tang2022}. Figure Q3-1 summarizes these reported ranges. They are
\textbf{context values, not universal pass/fail thresholds}.\par}

\begin{minipage}[t]{0.55\textwidth}
  \vspace{0pt}
  \centering
  \begin{tikzpicture}
  \begin{axis}[
    width=0.90\linewidth,
    height=63mm,
    xmin=0,xmax=65,
    xlabel={Reported interface stress (kPa)},
    ytick={1,2,3,4},
    yticklabels={Tang $p$,Tang $\tau$,Sanders tibial $\tau$,Sanders posterior $\tau$},
    y dir=reverse,
    tick label style={font=\fontsize{7.6}{8.6}\selectfont},
    label style={font=\fontsize{8.2}{9.2}\selectfont},
    xmajorgrids,
    grid style={AUBLine},
    axis line style={AUBGray},
    every axis plot/.append style={thick}]
    \addplot[only marks,mark=*,mark size=2.4pt,color=AUBRed,
      error bars/.cd,x dir=both,x explicit]
      coordinates {
        (57,1) +- (2,0)
        (15.5,2) +- (3.5,0)
        (22.3,3) +- (16.7,0)
        (22.85,4) +- (17.85,0)
      };
  \end{axis}
  \end{tikzpicture}
\end{minipage}\hfill
\begin{minipage}[t]{0.42\textwidth}
  \vspace{0pt}
  \centering
  {\fontsize{8.4}{9.7}\selectfont\bfseries\color{AUBRed}
  Reported shear-peak timing during stance}\par
  \vspace{1.0em}
  \begin{tikzpicture}[x=0.075\linewidth,y=1cm]
    \draw[->,color=AUBGray] (0,0) -- (10.35,0);
    \foreach \x/\lab in {0/0,2.5/25,4/40,6.5/65,8.5/85,10/100}{
      \draw[color=AUBGray] (\x,0.08)--(\x,-0.08)
        node[below,font=\fontsize{7}{8}\selectfont] {\lab};}
    \fill[AUBPale] (2.5,0.18) rectangle (4,0.82);
    \draw[AUBRed,thick] (2.5,0.18) rectangle (4,0.82);
    \fill[AUBPale] (6.5,0.18) rectangle (8.5,0.82);
    \draw[AUBRed,thick] (6.5,0.18) rectangle (8.5,0.82);
    \node[font=\fontsize{7.3}{8.3}\selectfont] at (3.25,0.5) {Peak 1};
    \node[font=\fontsize{7.3}{8.3}\selectfont] at (7.5,0.5) {Peak 2};
    \node[below,font=\fontsize{7.3}{8.4}\selectfont] at (5,-0.55)
      {Stance phase (\%)};
  \end{tikzpicture}
  \vspace{1.1em}

  {\fontsize{8.2}{9.7}\selectfont\raggedright
  The Q3 dynamic protocol therefore contains repeated loading and unloading
  rather than only a static maximum. Cyclic loading also exposes
  hysteresis, priming, and recovery in the PDMS force cell
  \cite{albrecht2019}.\par}
\end{minipage}

\vspace{0.3em}
{\fontsize{8.25}{9.6}\selectfont\color{AUBGray}
\textbf{Figure Q3-1. Literature-derived mechanical context.} Left: reported
min--max ranges from Tang \textit{et al.} and Sanders \textit{et al.}; right:
reported stance-phase windows for the two common shear peaks
\cite{q3sanders1992,q3tang2022}. Values describe their participants, sensor
locations, and protocols and are not transferred as clinical thresholds.\par}

\vspace{0.6em}
{\fontsize{9.4}{11.1}\selectfont
Zhang and Roberts reported a 226 kPa finite-element peak pressure at the
patellar tendon and a 50 kPa peak shear stress at the anterolateral tibia in
their participant-specific model and measurements \cite{q3zhang2001}. The
large difference between reported studies reinforces the need to record
anatomical location, activity, contact area, and duration and to avoid a single
global pressure threshold. The 2026 systematic review by Armitage
\textit{et al.} similarly found no consistent pressure threshold for comfort or
discomfort across regions and activities \cite{q3armitage2026}.\par}

\clearpage
\sectiontitle{Q3C. Thermal and Moisture Literature}

\hypertarget{q3-microclimate-evidence}{}
\qtwosubsection{3.3 Expected time-dependent microclimate behaviour}

{\fontsize{9.5}{11.25}\selectfont
Patern\`o \textit{et al.} instrumented a personalized transtibial liner with
12 temperature/RH loggers and used a protocol comprising 50 min seated rest,
30 min treadmill walking, and 50 min recovery \cite{paterno2020}. Temperature
rose during activity and decreased after activity, while RH continued to rise.
Figure Q3-2 replots their tabulated means and standard deviations. This pattern
supports time-series assessment rather than reliance on a single end-point.\par}

\begin{minipage}[t]{0.49\textwidth}
  \vspace{0pt}\centering
  \begin{tikzpicture}
  \begin{axis}[
    width=\linewidth,height=64mm,
    xmin=0.75,xmax=4.25,ymin=29.5,ymax=37,
    ylabel={Temperature ($^{\circ}$C)},
    xtick={1,2,3,4},
    xticklabels={R1 start,R1 end,Walk end,R2 end},
    tick label style={font=\fontsize{7.3}{8.3}\selectfont},
    xticklabel style={font=\fontsize{6.7}{7.6}\selectfont},
    label style={font=\fontsize{8.1}{9.1}\selectfont},
    legend style={font=\fontsize{7}{7.9}\selectfont,draw=none,at={(0.5,1.02)},anchor=south,legend columns=2},
    ymajorgrids,grid style={AUBLine},axis line style={AUBGray}]
    \addplot[color=AUBRed,mark=*,thick,error bars/.cd,y dir=both,y explicit]
      coordinates {(1,31.42) +- (0,0.86) (2,31.20) +- (0,1.01)
                   (3,34.62) +- (0,0.57) (4,33.91) +- (0,0.63)};
    \addlegendentry{Personalized liner}
    \addplot[color=AUBGray,mark=square*,thick,error bars/.cd,y dir=both,y explicit]
      coordinates {(1,32.91) +- (0,0.32) (2,33.52) +- (0,0.38)
                   (3,35.72) +- (0,0.28) (4,34.80) +- (0,0.49)};
    \addlegendentry{Pe-lite system}
  \end{axis}
  \end{tikzpicture}
\end{minipage}\hfill
\begin{minipage}[t]{0.49\textwidth}
  \vspace{0pt}\centering
  \begin{tikzpicture}
  \begin{axis}[
    width=\linewidth,height=64mm,
    xmin=0.75,xmax=4.25,ymin=60,ymax=89,
    ylabel={Relative humidity (\%RH)},
    xtick={1,2,3,4},
    xticklabels={R1 start,R1 end,Walk end,R2 end},
    tick label style={font=\fontsize{7.3}{8.3}\selectfont},
    xticklabel style={font=\fontsize{6.7}{7.6}\selectfont},
    label style={font=\fontsize{8.1}{9.1}\selectfont},
    legend style={font=\fontsize{7}{7.9}\selectfont,draw=none,at={(0.5,1.02)},anchor=south,legend columns=2},
    ymajorgrids,grid style={AUBLine},axis line style={AUBGray}]
    \addplot[color=AUBRed,mark=*,thick,error bars/.cd,y dir=both,y explicit]
      coordinates {(1,68.13) +- (0,1.33) (2,77.37) +- (0,2.11)
                   (3,79.87) +- (0,2.39) (4,83.48) +- (0,2.28)};
    \addlegendentry{Personalized liner}
    \addplot[color=AUBGray,mark=square*,thick,error bars/.cd,y dir=both,y explicit]
      coordinates {(1,71.25) +- (0,3.54) (2,77.94) +- (0,3.48)
                   (3,80.37) +- (0,3.01) (4,83.07) +- (0,3.34)};
    \addlegendentry{Pe-lite system}
  \end{axis}
  \end{tikzpicture}
\end{minipage}

\vspace{0.25em}
{\fontsize{8.2}{9.55}\selectfont\color{AUBGray}
\textbf{Figure Q3-2. Published temperature and RH progression.} Replotted from
Tables 2 and 3 of Patern\`o \textit{et al.}; points are the reported mean across
12 sensors and error bars are the reported standard deviation
\cite{paterno2020}. R1 start: initial rest; R1 end: end of 50-min rest;
Walk end: end of 30-min walking; R2 end: end of 50-min recovery. This was a
single-participant case study and does not define universal thresholds.\par}

\vspace{0.55em}
\begingroup
\fontsize{8.7}{10.3}\selectfont
\renewcommand{\arraystretch}{1.12}
\begin{tabularx}{\linewidth}{>{\raggedright\arraybackslash\bfseries\color{AUBRed}}p{3.0cm}Y Y}
\rowcolor{AUBRed}
\color{white}\textbf{Literature observation} &
\color{white}\textbf{Implication for Q3} &
\color{white}\textbf{Necessary control} \\
Temperature plateaued after donning and rose rapidly during walking
\cite{paterno2020}. & Include equilibration, a controlled heat step or ramp,
and a recovery interval. & Reference thermocouple/RTD at the same contact
location; record ambient temperature. \\
\rowcolor{AUBPale}
RH continued rising through activity and recovery \cite{paterno2020}. & Do not
expect humidity response to follow the temperature waveform or recover at the
same rate. & Reference hygrometer or calibrated chamber; document airflow and
sealed volume. \\
Multimodal pressure, temperature, and moisture surface with independent event observations
\cite{q3woo2022}. & Validate event labels against external observations rather
than declaring success from the sensor output itself. & Blinded event log or
automated fixture log containing true event time and location. \\
\rowcolor{AUBPale}
Soft prosthetic-interface platform with concurrent pressure and temperature
\cite{q3kwak2020}. & Time-synchronize modalities and
retain separate physical-unit traces before any combined display. & Common
clock, timestamp audit, and packet-loss report. \\
\end{tabularx}
\endgroup
\sourceanchor{Literature-derived observations and project test controls from
prosthetic-liner and multimodal interface-monitoring studies
\cite{paterno2020,q3woo2022,q3kwak2020}.}

\clearpage
\sectiontitle{Q3C. Evidence Architecture and Ground Truth}

\hypertarget{q3-ground-truth}{}
\qtwosubsection{3.4 Specimens, fixtures, and ground truth}

{\fontsize{9.55}{11.3}\selectfont
The primary test article is the directly printed 56 mm $\times$ 55 mm
$\times$ 5.2 mm PDMS slab defined in Q2. The PET print remains a fabrication
control where a comparable electrical measurement is possible; it is not a
liner-material alternative. At least three independently fabricated PDMS
specimens and five repeated cycles per condition define the replication target.
The specimen register records the achieved count and keeps independent slabs
distinct from repeated cycles.
Every file records specimen ID, channel ID, CAD coordinate, calibration
version, stimulus program, reference-instrument ID, sampling rate, and ambient
conditions.\par}

\begin{center}
\begin{tikzpicture}[node distance=7mm,>=stealth]
  \tikzstyle{qbox}=[draw=AUBLine,fill=AUBPale,rounded corners=1.5pt,
    minimum height=11mm,text width=30mm,align=center,
    font=\fontsize{8}{9}\selectfont]
  \node[qbox] (stim) {Programmed stimulus\\force, motion, heat, RH};
  \node[qbox,right=of stim] (slab) {Printed PDMS slab\\known CAD coordinate};
  \node[qbox,right=of slab] (daq) {Synchronized acquisition\\raw sensor channels};
  \node[qbox,right=of daq] (eval) {Blind comparison\\signal vs. ground truth};
  \draw[->,thick,color=AUBRed] (stim)--(slab);
  \draw[->,thick,color=AUBRed] (slab)--(daq);
  \draw[->,thick,color=AUBRed] (daq)--(eval);
  \draw[->,thick,dashed,color=AUBGray]
    ([yshift=-4mm]stim.south east) to[bend right=18]
    node[below,font=\fontsize{7}{8}\selectfont]{reference instrument and fixture log}
    ([yshift=-4mm]eval.south west);
\end{tikzpicture}
\end{center}

{\fontsize{8.2}{9.5}\selectfont\color{AUBGray}
\textbf{Figure Q3-3. Ground-truth validation architecture.} The slab output is
not used to label its own test. Stimulus magnitude, time, direction, and
coordinate are measured independently and then compared with the synchronized
sensor record.\par}

\vspace{0.55em}
\begingroup
\fontsize{8.1}{9.45}\selectfont
\setlength{\tabcolsep}{2.8pt}
\renewcommand{\arraystretch}{1.11}
\begin{tabularx}{\linewidth}{>{\centering\arraybackslash\bfseries\color{AUBRed}}p{0.75cm}
                            >{\raggedright\arraybackslash\bfseries}p{2.65cm}
                            >{\raggedright\arraybackslash}p{4.35cm}
                            >{\raggedright\arraybackslash}p{4.0cm}Y}
\rowcolor{AUBRed}
\color{white}\textbf{ID} & \color{white}\textbf{Condition} &
\color{white}\textbf{Programmed input} &
\color{white}\textbf{Independent reference} &
\color{white}\textbf{Primary evidence} \\
B0 & Baseline & Unloaded slab after thermal/RH equilibration; electronics on
and fixture untouched. & Reference temperature/RH; zero-force check; logged
DAQ offset. & Noise, zero drift, channel covariance, false-event rate. \\
\rowcolor{AUBPale}
P1 & Local pressure & Normal-force staircase and cyclic load at sensor centres,
between sensors, and near slab boundaries; fixed contactor area. & Calibrated
load cell or force gauge plus measured contact area. & Pressure sensitivity,
linearity, hysteresis, recovery, coordinate response. \\
S1 & Directional shear & Constant normal preload followed by signed tangential
force in both slab axes. & Two-axis force reference or tangential force gauge;
stage direction log. & Shear magnitude/direction, cross-axis error, pressure
cross-talk. \\
\rowcolor{AUBPale}
SL1 & Slip/pistoning & Programmed translation under contact, initially spanning
0.5--4.5 mm as a literature-anchored pilot range \cite{q3yu2024}. & Encoder,
LVDT/digital indicator, or calibrated video displacement. & Event detection,
direction, onset-time error, and displacement-linked spatial change. \\
T1 & Thermal buildup & Equilibrated baseline followed by 25--45 $^{\circ}$C
steps or a body-relevant subset; localized and uniform heating. & Contact
thermocouple or traceable RTD at the same coordinate. & Temperature RMSE,
response/recovery time, spatial leakage, force sensitivity. \\
\rowcolor{AUBPale}
H1 & Vapour humidity & 20--95\%RH chamber steps or the available validated
subset, following the calibration range used by Patern\`o \textit{et al.}
\cite{paterno2020}. & Calibrated chamber or reference hygrometer positioned
beside the slab. & RH RMSE, response/recovery, saturation, temperature coupling. \\
W1 & Liquid sweat & Known artificial-sweat dose normalized by exposed area;
separate deposition, dwell, and drying stages. & Calibrated pipette or gravimetric
mass; image/time log of wetted area. & Detection, saturation, leakage current,
recovery, permanent drift, damage. \\
\rowcolor{AUBPale}
C1 & Combined events & Pressure+shear, temperature+RH, pressure+liquid, then all
available modalities at one or separated coordinates. & All corresponding
references synchronized to one clock. & Missed events, false events,
cross-sensitivity, and degradation from isolated-condition performance. \\
\end{tabularx}
\endgroup
\sourceanchor{Implementation-ready test matrix. Reference methods and pilot ranges
are anchored to published force, motion, and liner microclimate studies
\cite{q3sanders1992,q3yu2024,paterno2020}; they are not clinical thresholds.}

\clearpage
\sectiontitle{Q3C. Published Phantom Precedents and Resources}

\hypertarget{q3-phantom-precedents}{}
\qtwosubsection{3.4.1 Three literature precedents retained as build options}

{\fontsize{9.3}{11}\selectfont
No single phantom reproduces pressure, shear, liquid sweat, vapour humidity,
temperature, and residual-limb volume change. The three candidates below are
therefore retained as complementary, published build options. The photographs
are reproduced from the cited papers, and every unchecked box means
\emph{availability at AUB Biomedical Engineering has not yet been confirmed}.\par}

\qtwosubsection{Candidate A --- Patern\`o variable-volume transfemoral simulator}

\noindent
\begin{minipage}[t]{0.385\linewidth}
\vspace{0pt}
\centering
\includegraphics[width=\linewidth]{assets/q3/phantoms/paterno2025_fig5a.png}
\vspace{0.2em}

{\fontsize{7.45}{8.7}\selectfont\raggedright\color{AUBGray}
\textbf{Figure Q3-4A. Published transfemoral simulator.} Cropped from
Patern\`o \textit{et al.}, Figure~5a: anatomical hip model, silicone tissue
layers, hydraulic chamber, and two syringe pumps \cite{paterno2025socket}.
CC BY 4.0; only the panel boundary was cropped.\par}
\end{minipage}\hfill
\begin{minipage}[t]{0.59\linewidth}
\vspace{0pt}
{\fontsize{8.05}{9.35}\selectfont
\textbf{Best-matched role.} Controlled whole-limb volume/preload change and
pressure-map response in transfemoral geometry. The published chamber produced
reversible changes up to approximately $\pm7.5\%$, with volume checked by 3D
scanning \cite{paterno2025socket}.\par}
\vspace{0.15em}
\begingroup
\fontsize{7.55}{8.75}\selectfont
\setlength{\tabcolsep}{2.3pt}
\renewcommand{\arraystretch}{1.08}
\begin{tabularx}{\linewidth}{>{\bfseries\color{AUBRed}}p{2.05cm}Y
 >{\centering\arraybackslash}p{1.05cm}}
\rowcolor{AUBRed}
\color{white}\textbf{Subsystem} & \color{white}\textbf{Required items} &
\color{white}\textbf{AUB?} \\
Geometry & Functional hip-joint model or equivalent; 3D-printed proximal
attachment; rigid mounting frame. & $\square$ \\
\rowcolor{AUBPale}
Tissue body & Molds and graded silicone layers representing muscle, fat, and
skin; release agent and casting tools. & $\square$ \\
Hydraulics & Sealed internal water chamber, tubing, fittings, valves, and
leak-test supplies. & $\square$ \\
\rowcolor{AUBPale}
Actuation & Two motorized syringe pumps (published: NE-1010) or validated
equivalent; syringes and water reservoir. & $\square$ \\
Ground truth & 3D scanner or validated geometric-volume method; reference
pressure and displacement instruments. & $\square$ \\
\end{tabularx}
\endgroup
\end{minipage}

\vspace{0.45em}
\qtwosubsection{Candidate B --- McGrath sweating residuum/socket simulator}

\noindent
\begin{minipage}[t]{0.385\linewidth}
\vspace{0pt}
\centering
\includegraphics[width=\linewidth]{assets/q3/phantoms/mcgrath2021_fig1.jpg}
\vspace{0.2em}

{\fontsize{7.45}{8.7}\selectfont\raggedright\color{AUBGray}
\textbf{Figure Q3-4B. Published sweating-interface simulator.} McGrath
\textit{et al.}, Figure~1: printed bones, porous silicone residuum, water inlet,
liner/check socket, and cyclic test-machine mounting
\cite{mcgrath2021sweat}. CC BY 4.0.\par}
\end{minipage}\hfill
\begin{minipage}[t]{0.59\linewidth}
\vspace{0pt}
{\fontsize{8.05}{9.35}\selectfont
\textbf{Best-matched role.} Liquid-moisture and moisture-induced axial slip
testing. The paper applied 20~mL of water through molded pores and compared dry
and wet cyclic displacement with standard and perforated liners
\cite{mcgrath2021sweat}. This is not a vapour-RH calibration chamber.\par}
\vspace{0.15em}
\begingroup
\fontsize{7.55}{8.75}\selectfont
\setlength{\tabcolsep}{2.3pt}
\renewcommand{\arraystretch}{1.08}
\begin{tabularx}{\linewidth}{>{\bfseries\color{AUBRed}}p{2.05cm}Y
 >{\centering\arraybackslash}p{1.05cm}}
\rowcolor{AUBRed}
\color{white}\textbf{Subsystem} & \color{white}\textbf{Required items} &
\color{white}\textbf{AUB?} \\
Geometry & Cast/scan-derived transtibial shape; printed bone model; two-part
negative mold. & $\square$ \\
\rowcolor{AUBPale}
Porous body & Silicone soft-tissue body; 3-mm pore-forming straws/rods; silicone
adhesive and seals. & $\square$ \\
Water delivery & Graduated syringe, rubber tubing, proximal inlet, measured
water dose, and drying arrangement. & $\square$ \\
\rowcolor{AUBPale}
Socket interface & Pin-lock check socket, lock, standard liner, and perforated
comparison liner. & $\square$ \\
Cyclic reference & Universal test machine or controlled axial load fixture;
force and displacement logging. & $\square$ \\
\end{tabularx}
\endgroup
\end{minipage}

\clearpage
\sectiontitle{Q3C. Phantom Selection Logic}

\qtwosubsection{Candidate C --- Phillips dual-hardness, water-bladder mock limb}

\noindent
\begin{minipage}[t]{0.56\linewidth}
\vspace{0pt}
\centering
\includegraphics[width=\linewidth]{assets/q3/phantoms/phillips2025_fig1.png}
\vspace{0.2em}

{\fontsize{7.45}{8.7}\selectfont\raggedright\color{AUBGray}
\textbf{Figure Q3-4C. Published low-cost adjustable mock limb.} Phillips
\textit{et al.}, Figure~1: urethane casting sequence, TPU bladder, sensor
placement, rigid socket, Arduino, and motorized syringe system
\cite{phillips2025}. CC BY 4.0; only page margins were cropped.\par}
\end{minipage}\hfill
\begin{minipage}[t]{0.415\linewidth}
\vspace{0pt}
{\fontsize{8.0}{9.3}\selectfont
\textbf{Best-matched role.} A reproducible, relatively accessible adjustable
limb for pressure/preload and volume-change studies. The published build used
three water-filled bladders, reached controlled changes up to $\pm5\%$, and cost
less than CAD~400 \cite{phillips2025}.\par}
\vspace{0.15em}
\begingroup
\fontsize{7.35}{8.55}\selectfont
\setlength{\tabcolsep}{2.15pt}
\renewcommand{\arraystretch}{1.07}
\begin{tabularx}{\linewidth}{>{\bfseries\color{AUBRed}}p{1.7cm}Y
 >{\centering\arraybackslash}p{0.9cm}}
\rowcolor{AUBRed}
\color{white}\textbf{Subsystem} & \color{white}\textbf{Required items} &
\color{white}\textbf{AUB?} \\
Core/body & 30-mm aluminum mandrel; VytaFlex 60 core and VytaFlex 20 outer
layer or characterized equivalents. & $\square$ \\
\rowcolor{AUBPale}
Molds & Printed PLA inner, outer, lid, and bladder-cavity molds; epoxy seal,
release agent, and clamp stand. & $\square$ \\
Bladders & Three heat-sealed 6$\times$10-cm TPU bladders; 6.4-mm fittings,
tubing, Parafilm/seals, and water. & $\square$ \\
\rowcolor{AUBPale}
Actuation & Motorized syringe system, linear actuator, Arduino, motor driver,
power supply, and computer. & $\square$ \\
Interface / ref. & 5-mm PLA socket; reference pressure film/system;
volume and circumference measurements. & $\square$ \\
\end{tabularx}
\endgroup
\end{minipage}

\vspace{0.55em}
\qtwosubsection{3.4.2 Evidence-based selection for the project tests}

\begingroup
\fontsize{7.85}{9.15}\selectfont
\setlength{\tabcolsep}{2.5pt}
\renewcommand{\arraystretch}{1.09}
\begin{tabularx}{\linewidth}{>{\bfseries\color{AUBRed}}p{2.25cm}Y Y Y}
\rowcolor{AUBRed}
\color{white}\textbf{Published option} & \color{white}\textbf{Use it for} &
\color{white}\textbf{Evidence already demonstrated} &
\color{white}\textbf{Boundary requiring a separate fixture} \\
Patern\`o 2025 & Transfemoral volume/preload steps and sensor-map response to
global fit change. & Hydraulic chamber, two pumps, repeated 3D volume
measurement, and approximately $\pm7.5\%$ range. & Does not generate liquid
sweat or controlled vapour RH; comparatively complex fabrication. \\
\rowcolor{AUBPale}
McGrath 2021 & Liquid-sweat exposure and wet-interface slip under axial cyclic
loading. & Porous silicone body, measured 20-mL water dose, dry/wet liner
comparison, and displacement reference. & Transtibial geometry; no controlled
vapour RH or temperature field; liquid distribution was not perfectly uniform. \\
Phillips 2025 & Low-cost volume/preload and localized pressure changes using
independent bladders. & Dual-hardness body, three water bladders, Arduino
actuation, $\pm5\%$ change, and build cost below CAD~400. & No sweating pores or
RH control; published pressure film exhibited long-duration drift. \\
\end{tabularx}
\endgroup
\sourceanchor{The three options are published implementations, not speculative
phantom sketches \cite{paterno2025socket,mcgrath2021sweat,phillips2025}. Their
functions are complementary: volume/preload, liquid-moisture slip, and low-cost
localized adjustability, respectively.}

\vspace{0.35em}
\evidencebox{Implementation rule:}{Retain all three as possible constructions,
then select the first build through the AUB inventory register. Use the
Patern\`o or Phillips architecture for pressure/volume localization and the
McGrath architecture for liquid-sweat/slip. Vapour humidity and contact
temperature still require the separate calibrated chamber and reference probes
defined in Test Matrix H1/T1; no phantom substitutes for those ground truths.}

\vspace{0.35em}
\evidencebox{Safety and scope:}{Use guards, hard stops, and leak checks; begin
with small loads and place no person in the fixture. These phantoms validate
sensor response, coordinate mapping, and controlled interface conditions. They
do not establish human comfort, tissue equivalence, or clinical safety.}

\clearpage
\sectiontitle{Q3C. Measurement Models}

\hypertarget{q3-measurement-models}{}
\qtwosubsection{3.5 Converting raw channels into testable measurements}

{\fontsize{9.4}{11.1}\selectfont
All comparisons use baseline-corrected raw measurements. For channel $j$
of slab $i$, the baseline-corrected signal is
\begin{equation}
  \widetilde{s}_{ij}(t)=s_{ij}(t)-\overline{s}_{ij,0},
  \qquad
  \overline{s}_{ij,0}=\frac{1}{N_0}\sum_{k=1}^{N_0}s_{ij}(t_k),
  \label{eq:q3baseline}
\end{equation}
where the baseline window $N_0$ is recorded before the stimulus. Filtering may
remove electrical noise, while both raw and filtered data are retained and
the filter type, cutoff, phase delay, and missing samples reported. Thus,
Equation~\eqref{eq:q3baseline} makes the zeroing operation explicit and
reproducible.\par

\textbf{Pressure and shear.} The four-capacitor force cell is calibrated as a
multivariable system because normal and tangential loads can affect more than
one capacitor. For the calibration vector
$\Delta\mathbf c=[\Delta C_1,\Delta C_2,\Delta C_3,\Delta C_4]^T$,
\begin{equation}
  \Delta\mathbf c=\mathbf K
  \begin{bmatrix}F_x&F_y&F_z\end{bmatrix}^{T}+\boldsymbol\varepsilon,
  \qquad
  \widehat{\mathbf F}=\mathbf K^{+}\Delta\mathbf c,
  \label{eq:q3forcecal}
\end{equation}
where $\mathbf K$ is fitted from signed, known loads and $\mathbf K^{+}$ is its
pseudoinverse. This empirical matrix accommodates unequal printed channels and
is safer than assuming perfect geometric symmetry. Equation~\eqref{eq:q3forcecal}
is therefore the primary force-recovery model. The engineering stresses
under a documented contact area $A$ are
\begin{equation}
  p=\frac{F_z}{A},\qquad
  \tau_x=\frac{F_x}{A},\qquad
  \tau_y=\frac{F_y}{A},\qquad
  \tau=\sqrt{\tau_x^2+\tau_y^2}.
  \label{eq:q3stress}
\end{equation}
Albrecht \textit{et al.} found linear response between 0.1 and 8 N in their
printed cell, but also reported priming, drift, and shear hysteresis linked to
the viscoelastic PDMS spacer \cite{albrecht2019}; the project's own
$\mathbf K$ and working range are therefore obtained from the present slab
before applying
Equation~\eqref{eq:q3stress}.\par

\textbf{Temperature and humidity.} Calibration coefficients are fitted against
reference instruments rather than copied from the source papers:
\begin{equation}
 \widehat{T}=a_0+a_1R_T+a_2R_T^2,
 \qquad
 \widehat{RH}=b_0+b_1C_H+b_2C_H^2+b_3\widehat{T},
 \label{eq:q3microcal}
\end{equation}
where $R_T$ is the printed meander resistance and $C_H$ is the humidity-IDE
capacitance. Equation~\eqref{eq:q3microcal} allows curvature and thermal
compensation to be tested rather than presumed. The temperature term in the RH equation is retained only if
model comparison shows that it reduces reference error; otherwise it is
removed to avoid unnecessary fitting. Calibration, cross-validation, and Q3
testing data are held separate.\par

\textbf{Exposure over time.} A peak alone cannot describe a sustained event.
For a window $[t_0,t_1]$, report the pressure-time and shear-time integrals
\begin{equation}
 PTI=\int_{t_0}^{t_1}\max[p(t)-p_0,0]\,\mathrm{d}t,
 \qquad
 STI=\int_{t_0}^{t_1}\max[\tau(t)-\tau_0,0]\,\mathrm{d}t,
 \label{eq:q3dose}
\end{equation}
along with peak, mean, and dwell time. Here $p_0$ and $\tau_0$ are engineering
noise or event-detection levels established during B0, not clinical injury
thresholds; Equation~\eqref{eq:q3dose} therefore describes exposure above a
measured detection level, not biological damage.\par

\textbf{Slip label.} A trial is labelled as true controlled slip only when the
independent reference satisfies
\begin{equation}
 y_{\mathrm{slip}}(t)=
 \begin{cases}
 1,&|d_{\mathrm{ref}}(t)-d_{\mathrm{ref}}(t-\Delta t)|\ge d_{\min},\\
 0,&\text{otherwise},
 \end{cases}
 \label{eq:q3slip}
\end{equation}
where $d_{\min}$ is the verified resolution of the displacement fixture.
Sensor-based detection is then compared with the external label in
Equation~\eqref{eq:q3slip}. Without
$d_{\mathrm{ref}}$, the result is reported as a \emph{shear event}, not slip.\par}

\vspace{0.45em}
\begin{center}
\begin{tikzpicture}[>=stealth,node distance=2.7mm]
  \tikzstyle{calbox}=[draw=AUBLine,fill=AUBPale,rounded corners=1pt,
    minimum height=7mm,text width=31mm,align=center,
    font=\fontsize{6.9}{7.8}\selectfont]
  \node[calbox] (r) {raw channels +\\fixture references};
  \node[calbox,right=of r] (b) {baseline and\\quality audit};
  \node[calbox,right=of b] (m) {calibration matrix /\\microclimate model};
  \node[calbox,right=of m] (u) {physical units +\\uncertainty};
  \node[calbox,right=of u] (e) {event and\\location metrics};
  \draw[->,thick,color=AUBRed] (r)--(b);
  \draw[->,thick,color=AUBRed] (b)--(m);
  \draw[->,thick,color=AUBRed] (m)--(u);
  \draw[->,thick,color=AUBRed] (u)--(e);
\end{tikzpicture}
\end{center}
{\fontsize{7.5}{8.7}\selectfont\color{AUBGray}
\textbf{Figure Q3-5. Calibration-to-decision pipeline.} Raw channels remain
traceable to the external reference before any event label or map is produced.\par}
\sourceanchor{The equations are project calibration definitions. Sensor-specific
transfer functions are estimated from the present slab; published devices
provide mechanisms and candidate forms, not transferable coefficients
\cite{albrecht2019,liew2020,aitmammar2020}.}

\clearpage
\sectiontitle{Q3C. Condition-by-Condition Evidence Protocol}

\hypertarget{q3-execution-sequence}{}
\qtwosubsection{3.6 Defined execution sequence}

\begingroup
\fontsize{8.35}{9.8}\selectfont
\setlength{\tabcolsep}{3pt}
\renewcommand{\arraystretch}{1.13}
\begin{tabularx}{\linewidth}{>{\raggedright\arraybackslash\bfseries\color{AUBRed}}p{2.8cm}
                            >{\raggedright\arraybackslash}p{5.0cm}Y Y}
\rowcolor{AUBRed}
\color{white}\textbf{Protocol} &
\color{white}\textbf{Stimulus sequence} &
\color{white}\textbf{Control and negative test} &
\color{white}\textbf{Required plots} \\
P1: localized normal loading & Apply an unloaded baseline, increasing force
steps, dwell, decreasing steps, recovery, and repeated cyclic loading. Repeat
at every active footprint, at inter-sensor midpoints, and near slab edges. &
Repeat the fixture motion without contact; rotate or replace the contactor while
keeping force and area documented. & Raw and calibrated traces; applied versus
estimated force; residual plot; hysteresis loop; coordinate-response map. \\
\rowcolor{AUBPale}
S1: directional shear & Establish constant $F_z$, then apply $+x,-x,+y,-y$
force staircases and cycles. Include at least one priming cycle before the
reported cycles because PDMS viscoelasticity can cause first-cycle differences
\cite{albrecht2019}. & Normal preload without tangential motion; tangential
fixture motion without contact; repeat after rotating the slab 180 degrees. &
$\widehat F_x,\widehat F_y,\widehat F_z$ versus references; vector-angle error;
cross-axis matrix; loading/unloading curves. \\
SL1: controlled slip & Translate the contactor by known distances at several
speeds under low, medium, and high normal preload. Program forward, reverse,
stop-start, and no-slip high-shear trials. & High shear with no displacement is
the critical negative control; free motion without contact checks fixture
artefacts. Yu \textit{et al.} support direct displacement characterization
\cite{q3yu2024}. & Reference displacement and sensor traces on
one time axis; detected versus true event raster; onset-time and distance error. \\
\rowcolor{AUBPale}
T1: thermal buildup & Equilibrate, apply temperature steps or a controlled ramp,
hold, then remove heat and continue through recovery. Test localized contact
and whole-environment heating separately. & Reference sensor at the active
location; heater motion with heater off; force-only contact at constant
temperature. & Estimated versus reference temperature; RMSE; response and
recovery curves; spatial thermal spread. \\
H1: vapour humidity & Equilibrate at low RH, apply increasing chamber steps,
decreasing steps, and recovery. Repeat at two or more temperatures to expose
temperature coupling. & Sealed dry-air step; chamber program with unpowered
sensor; reference probe placed beside the active region. & Capacitance and
calibrated RH; reference residuals; hysteresis; response/recovery time. \\
\rowcolor{AUBPale}
W1: artificial sweat & Deposit a recorded areal dose, hold, remove surface
liquid, dry, and repeat after recovery. Inspect protection, conductors, leakage,
and permanent offset before the next dose. & Equal-volume deionized-water test
only if chemically relevant and approved; dry pipette contact; protected versus
unprotected route from Q2. & Wetting-area image; raw signal; saturation duration;
drying/recovery; pre/post calibration and leakage. \\
C1: combinations & Begin with pressure+shear and temperature+RH; progress to
events applied at the same coordinate, then separated coordinates, and finally
all available modalities. Randomize trial order. & Repeat each constituent as
an isolated event on the same specimen and day. & Synchronized small multiples;
isolated-versus-combined error; confusion matrix; missed and false event map. \\
\end{tabularx}
\endgroup
\sourceanchor{Execution protocol anchored to published printed
force-cell behaviour, direct displacement validation, and liner microclimate
measurements \cite{albrecht2019,noll2018,paterno2020}.}

\vspace{0.6em}
\evidencebox{Order of testing:}{Run B0, isolated P1/S1/T1/H1, and the protected
liquid test W1 before slip and combined-condition trials. This order identifies
individual calibration and damage before several sources are superimposed.
Randomize trials within each completed block and do not tune thresholds on the
same repetitions used for final reporting.}

\vspace{0.65em}
\begin{center}
\begin{tikzpicture}[>=stealth,node distance=4mm]
  \tikzstyle{tracebox}=[draw=AUBLine,fill=AUBPale,rounded corners=1pt,
    minimum height=11mm,text width=40mm,align=center,
    font=\fontsize{7.2}{8.2}\selectfont]
  \node[tracebox] (branch) {select Figure Q3-4\\validation branch};
  \node[tracebox,right=of branch] (truth) {record independent\\ground truth};
  \node[tracebox,right=of truth] (sync) {synchronize with\\raw slab channels};
  \node[tracebox,right=of sync] (plots) {produce Section 3.8\\reference plots};
  \draw[->,thick,color=AUBRed] (branch)--(truth);
  \draw[->,thick,color=AUBRed] (truth)--(sync);
  \draw[->,thick,color=AUBRed] (sync)--(plots);
\end{tikzpicture}
\end{center}
{\fontsize{7.6}{8.8}\selectfont\color{AUBGray}
\textbf{Figure Q3-6. Fixture-to-evidence trace.} Each validation branch records
a separately measured reference channel, synchronizes it with the raw slab
output, and carries both into the final plots.\par}

\clearpage
\sectiontitle{Q3C. Evidence Metrics and Decision Gates}

\hypertarget{q3-metrics}{}
\qtwosubsection{3.7 Measurement quality, detection, and localization}

{\fontsize{9.25}{10.95}\selectfont
For each condition, report the calibration model, full test range, number of
specimens, repetitions, mean, standard deviation, raw traces, and rejected
trials with reasons. The core accuracy and localization quantities are
\begin{align}
 RMSE_m &=\sqrt{\frac{1}{N}\sum_{k=1}^{N}
   \left(\widehat{x}_{m,k}-x_{m,k}^{\mathrm{ref}}\right)^2},
 & NRMSE_m&=\frac{RMSE_m}{x_{m,\max}^{\mathrm{ref}}-x_{m,\min}^{\mathrm{ref}}},
 \label{eq:q3rmse}\\
 \widehat{\mathbf r}(t)&=
 \frac{\sum_i w_i(t)\mathbf r_i}{\sum_i w_i(t)},
 & e_{\mathrm{loc}}&=\left\|\widehat{\mathbf r}-\mathbf r_{\mathrm{true}}\right\|_2,
 \label{eq:q3loc}
\end{align}
where $m$ denotes modality, $\mathbf r_i$ is the CAD coordinate of sensor $i$,
and $w_i$ is a non-negative baseline-corrected response for the modality under
test. Equation~\eqref{eq:q3rmse} evaluates reference agreement, while
Equation~\eqref{eq:q3loc} evaluates spatial recovery. The centroid is an
engineering localization estimator used for Q3; Q5 evaluates and selects the
final interpolation algorithm.\par

\begin{equation}
 H_m=\frac{\max_x|y_{\uparrow}(x)-y_{\downarrow}(x)|}{y_{\mathrm{FS}}},
 \qquad
 XT_{m\rightarrow n}=\frac{\max|\Delta y_n|}{\max|\Delta y_m|},
 \label{eq:q3hystxt}
\end{equation}
Equation~\eqref{eq:q3hystxt} defines full-scale-normalized hysteresis and
cross-sensitivity from stimulated
modality $m$ into non-target modality $n$. Response time $t_{90}$ is the time
to reach 90\% of the final step response; recovery time is defined analogously
after removal. For event labels, report sensitivity, specificity, precision,
and F1 score from the complete confusion matrix rather than accuracy alone.\par}

\begingroup
\fontsize{8.25}{9.7}\selectfont
\setlength{\tabcolsep}{3pt}
\renewcommand{\arraystretch}{1.12}
\begin{tabularx}{\linewidth}{>{\raggedright\arraybackslash\bfseries\color{AUBRed}}p{3.2cm}
                            >{\raggedright\arraybackslash}p{4.1cm}Y Y}
\rowcolor{AUBRed}
\color{white}\textbf{Gate} &
\color{white}\textbf{Provisional engineering target} &
\color{white}\textbf{Report required} &
\color{white}\textbf{Interpretation if failed} \\
Calibration fit & $R^2\ge0.95$ and NRMSE $\le10\%$ of calibrated full scale
for each modality. & Per-specimen fit and held-out-test error; residuals versus
input and time. & Revise geometry, electronics, range, or calibration model;
do not hide nonlinear regions. \\
\rowcolor{AUBPale}
Repeatability & Coefficient of variation $\le10\%$ at repeated levels after
the stated priming procedure. & Within-cycle, between-cycle, and
between-specimen variation. & Manufacturing or viscoelastic response is not
yet sufficiently controlled. \\
Hysteresis & $H_m\le10\%$ full scale over the declared operating range. & Full
loading/unloading loops at slow and dynamic rates. & Restrict the range or use
history-aware calibration; report the limitation. \\
\rowcolor{AUBPale}
Cross-sensitivity & Non-target response $\le15\%$ of the target full-scale
response in isolated tests. & Complete modality-by-modality and axis-by-axis
cross-talk matrix. & Independent interpretation requires redesign or
multivariate compensation when this gate is not met. \\
Event detection & Sensitivity and specificity each $\ge90\%$ on held-out
repetitions. & Confusion matrix with event counts, not percentages alone. &
Thresholds or features are not transferable to unseen repetitions. \\
\rowcolor{AUBPale}
Localization & Median error no greater than half the nearest-neighbour sensor
pitch; 95th percentile also reported. & Error by coordinate, boundary location,
and isolated/combined condition. & Sensor density or coordinate mapping is
insufficient; this result feeds Q6. \\
Combined-condition robustness & No more than 10 percentage-point loss in event
sensitivity relative to isolated tests. & Paired isolated/combined comparison
on the same specimens. & Cross-sensitivity or signal fusion masks a modality. \\
\rowcolor{AUBPale}
Damage and recovery & No open/short circuit, leakage, delamination, permanent
saturation, or greater than 10\% full-scale baseline shift after recovery. &
Electrical inspection and pre/post calibration for every specimen. & Stop
progression to combined tests or liner integration. \\
\end{tabularx}
\endgroup

\vspace{0.4em}
{\fontsize{8.25}{9.55}\selectfont\color{AUBGray}
\textbf{Table note.} These numerical gates are the selected engineering
targets for preregistration before data collection. They are not published
clinical safety or discomfort thresholds and may be revised only before the
held-out Q3 evaluation, with the revision documented.\par}
\sourceanchor{Preregistration gates selected for this development stage. They are not
presented as clinical limits; the literature review reports substantial
participant, location, and protocol variability \cite{q3armitage2026}.}

\vspace{0.35em}
\begin{center}
\begin{tikzpicture}[>=stealth,node distance=3mm]
  \tikzstyle{gatebox}=[draw=AUBLine,fill=AUBPale,rounded corners=1pt,
    minimum height=8mm,text width=39mm,align=center,
    font=\fontsize{7}{8}\selectfont]
  \node[gatebox] (g1) {1. isolated-modality\\calibration pass};
  \node[gatebox,right=of g1] (g2) {2. combined-event\\separation pass};
  \node[gatebox,right=of g2] (g3) {3. post-exposure\\integrity pass};
  \node[gatebox,right=of g3] (g4) {4. proceed to\\curved integration};
  \draw[->,thick,color=AUBRed] (g1)--(g2);
  \draw[->,thick,color=AUBRed] (g2)--(g3);
  \draw[->,thick,color=AUBRed] (g3)--(g4);
\end{tikzpicture}
\end{center}
{\fontsize{7.5}{8.7}\selectfont\color{AUBGray}
\textbf{Figure Q3-7. Sequential engineering gates.} Failure at an earlier
gate blocks the more complex claim rather than being averaged into a total
score.\par}

\clearpage
\sectiontitle{Q3C. Validation Outputs and Interpretation}

\hypertarget{q3-validation-outputs}{}
\qtwosubsection{3.8 Evidence output formats and interpretation criteria}

{\fontsize{9.25}{10.95}\selectfont
The assignment asks what remains to be learned or proven and how the prototype
is evaluated, including tools, rationale, expected outcome, and design impact.
The panels below consolidate the reporting outputs paired with the
project-specific proof relationships in the Q3C matrix. They define the
evidence format and are not presented as completed measurements.\par}

\begin{minipage}[t]{0.485\textwidth}
  \plannedoutput{42mm}{Calibration and transient plots}{
  Required contents: reference-versus-estimate plots with held-out points,
  raw and filtered traces, hysteresis, response/recovery markers, error bars,
  and specimen IDs for pressure, shear, temperature, and RH.}
\end{minipage}\hfill
\begin{minipage}[t]{0.485\textwidth}
  \plannedoutput{42mm}{Known versus detected event locations}{
  Required contents: CAD sensor coordinates, true stimulus coordinates,
  detected centroids, error vectors, sensor footprints, and a
  confidence/coverage mask. Do not interpolate outside validated coverage.}
\end{minipage}

\vspace{0.45em}
\begin{minipage}[t]{0.485\textwidth}
  \plannedoutput{39mm}{Slip ground-truth comparison}{
  Required contents: synchronized reference displacement, normal pressure, signed
  shear, detected event state, onset-time error, and forward/reverse trials.}
\end{minipage}\hfill
\begin{minipage}[t]{0.485\textwidth}
  \plannedoutput{39mm}{Combined-condition small multiples}{
  Required contents: separate physical-unit traces and maps before any
  combined score; show disagreement, false events, missed events, and
  isolated-condition baselines.}
\end{minipage}

\vspace{0.5em}
\begingroup
\fontsize{8.55}{10.05}\selectfont
\setlength{\tabcolsep}{3pt}
\renewcommand{\arraystretch}{1.14}
\begin{tabularx}{\linewidth}{>{\raggedright\arraybackslash\bfseries\color{AUBRed}}p{3.4cm}Y Y}
\rowcolor{AUBRed}
\color{white}\textbf{Evidence field} &
\color{white}\textbf{Recorded evidence} &
\color{white}\textbf{Design interpretation} \\
Test population & Number of independently fabricated slabs, cycles, stimulus
coordinates, and accepted/rejected trials. & Do not treat repeated cycles on
one slab as independent specimens. \\
\rowcolor{AUBPale}
Reference agreement & $R^2$, RMSE/NRMSE, bias, 95\% limits of agreement where
sample size supports them, and residual plots. & A high $R^2$ does not remove
systematic bias; always report an error metric. \\
Dynamic response & Peak, mean, integral, $t_{90}$, recovery, hysteresis, and
drift for each modality. & Compare with the condition's reference timescale;
do not require humidity to respond as fast as pressure. \\
\rowcolor{AUBPale}
Detection performance & Full confusion matrix, sensitivity, specificity,
precision, and F1 by event type. & Results are engineering event recognition,
not diagnostic sensitivity for injury or discomfort. \\
Spatial performance & Median, interquartile range, and 95th-percentile
localization error; boundary and inter-sensor errors shown separately. & A
visually smooth map is not evidence of correct localization. \\
\rowcolor{AUBPale}
Combined-condition effect & Difference between isolated and combined trials in
error, sensitivity, response time, and baseline. & If modalities disagree,
retain separate outputs; do not average away the disagreement. \\
Post-exposure integrity & Continuity, insulation/leakage, visible damage,
baseline and sensitivity changes after liquid, heat, shear, and cyclic load. &
Passing detection before damage does not establish reusability. \\
\end{tabularx}
\endgroup
\sourceanchor{Evidence reporting template aligned with the external
reference architecture in Figure Q3-3 and the published distinction between
force and relative motion \cite{q3tang2022,noll2018}.}

\vspace{0.65em}
\evidencebox{Connection to Q5 software work:}{Q3 uses the CAD
coordinates only to verify whether a known event was detected at the correct
location. Detailed interpolation, threshold fusion, real-time rendering, and
the selected interface-condition heatmap architecture are developed in Q5.
Turner \textit{et al.} provide a key reference for comparing radial-basis,
Clough--Tocher, nearest-neighbour, and linear interpolation on a prosthetic
socket model \cite{q3turner2021}.}

\clearpage
\sectiontitle{Q3C. F = Interim Functional-Block Verdict}

\begin{center}
\begin{tikzpicture}[>=stealth,node distance=2.7mm]
  \tikzstyle{maturity}=[draw=AUBLine,fill=AUBPale,rounded corners=1pt,
    minimum height=8mm,text width=36mm,align=center,
    font=\fontsize{7}{8}\selectfont]
  \node[maturity] (m1) {literature\\plausibility};
  \node[maturity,right=of m1] (m2) {flat-slab\\calibration};
  \node[maturity,right=of m2] (m3) {phantom\\condition detection};
  \node[maturity,right=of m3] (m4) {subsequent participant-\\specific validation};
  \draw[->,thick,color=AUBRed] (m1)--(m2);
  \draw[->,thick,color=AUBRed] (m2)--(m3);
  \draw[->,thick,dashed,color=AUBGray] (m3)--(m4);
\end{tikzpicture}
\end{center}
{\fontsize{7.6}{8.8}\selectfont\color{AUBGray}
\textbf{Figure Q3-8. Claim maturity at the present project stage.} Q3 ends at
controlled phantom detection; the dashed clinical step is outside the current
evidence.\par}

\begingroup
\setlength{\fboxsep}{7pt}
\fcolorbox{AUBRed}{AUBPale}{%
\begin{minipage}{0.96\linewidth}
{\fontsize{12.2}{14.2}\selectfont\bfseries\color{AUBRed}
F = FUNCTIONAL BLOCKS DEFINED; LITERATURE BASIS ESTABLISHED}
\par
\vspace{0.35em}
{\fontsize{9.45}{11.15}\selectfont
\begin{itemize}[leftmargin=1.25em,itemsep=0.3em,topsep=0.1em]
  \item \textbf{Leveraged knowledge:} prior studies support the underlying
  mechanical, thermal, humidity, printed-PDMS, relative-motion, and spatial
  visualization mechanisms
  \cite{q3sanders1992,q3tang2022,paterno2020,q3kwak2020,q3armitage2026}.
  \item \textbf{Prototype-specific proof:} the right-hand column of the Q3C
  matrix identifies the relationships that are not transferable from those
  studies: calibration, separation, localization, recovery, and integrity for
  this slab and sensor placement.
  \item \textbf{Decision rule:} report ``yes,'' ``partial,'' or ``no'' from the
  preregistered gates in Sections 3.6--3.8. A partial result names the
  passing and failing modalities; an overall average cannot conceal failed
  localization, separation, recovery, or integrity.
  \item \textbf{Clinical boundary:} even a complete engineering pass supports
  only detection of controlled interface conditions. The words
  \emph{discomfort}, \emph{harmful}, or \emph{skin risk} require subsequent
  participant-specific validation and clinical oversight.
  \item \textbf{Evidence sequence:} B0 and isolated-condition calibration
  precede blinded slip, combined-condition, heatmap, and whole-liner claims.
\end{itemize}
\textbf{Current verdict: Q3A identifies the eight critical blocks, Q3B defines
their engineering boundaries, and Q3C distinguishes leveraged literature from
the specimen-level evidence required for the present prototype.}}
\end{minipage}}
\endgroup

\vspace{0.75em}
\evidencebox{Claim-evidence statement:}{This is the appropriate boundary for an
open-ended project progress report: the functional architecture, literature
basis, candidate fixtures, ground-truth comparisons, metrics, and evidence
formats are established. A claim that the slab \textbf{successfully detects}
the four conditions is reserved for specimen-level experimental results.}

\clearpage
